% chapters/14/evaluasi.tex
\bagian{Asesmen dan Instrumen Penilaian}
\label{sec:asesmen-aset}

\begin{asesmenbox}
\textbf{Kedudukan asesmen bab ini.} Bab ini menuntaskan \SUBCPMK{4.1}, yaitu
bab pertama pada \CPMK{4}. Ketiga indikatornya diperiksa melalui penilaian
artefak \istilah{project} (Tugas Praktikum 11) dan soal konsep.

\vspace{0.4em}
Materi Bab~XIV termasuk bahan \textbf{Proyek Akhir} pada minggu kelima belas
dan \textbf{Ujian Akhir Semester} pada minggu keenam belas.
\end{asesmenbox}

\subsection*{Instrumen 1: Rubrik Praktikum 11}

\begin{table}[htbp]
\centering
\caption{Rubrik penilaian Praktikum 11 beserta indikator yang diukur}
\label{tab:rubrik-praktikum-11}
{\setlength{\tabcolsep}{4pt}\footnotesize%
\begin{tabularx}{\textwidth}{@{}L{3.2cm} C{1.3cm} C{2.1cm} Y@{}}
\toprule
\headrow \textbf{Kriteria} & \textbf{Bobot} & \textbf{Indikator} & \textbf{Deskripsi capaian} \\
\midrule
Impor dan pengaturan aset & 35\% & \ind{a} &
Aset terimpor dengan pengaturan yang tepat dan seragam; ukuran gambar pada
scene sesuai hitungan; sumber serta lisensi setiap aset tercatat. \\
\addlinespace
Animasi & 35\% & \ind{b} &
\istilah{Animator} memuat keadaan diam dan berlari beserta transisinya, dan
animasinya berganti dengan benar menurut gerak pemain. \\
\addlinespace
Antarmuka pengguna & 30\% & \ind{c} &
\istilah{Canvas} menampilkan skor dan nyawa yang benar-benar berubah mengikuti
keadaan permainan, dan keadaan permainannya tetap disimpan pada kode. \\
\bottomrule
\end{tabularx}}
\end{table}

\subsection*{Instrumen 2: Soal Konsep}

Jawablah dengan singkat. Kunci jawaban beserta indikator yang diukur disajikan
pada Tabel~\ref{tab:kunci-konsep-ch14}.

\begin{enumerate}
  \item Sebuah gambar berukuran 192 kali 192 piksel diimpor dengan
        \istilah{Pixels Per Unit} bernilai 100. Hitunglah tingginya pada scene,
        lalu hitunglah nilai yang diperlukan agar tingginya menjadi 2 satuan.
  \item Jelaskan mengapa ukuran berkas gambar tidak boleh diubah di luar Unity.
  \item Kode menetapkan nilai parameter animasi, tetapi animasinya tidak
        berganti dan tidak ada pesan kesalahan. Sebutkan dua kemungkinan
        penyebabnya.
  \item Parameter jenis apa yang sebaiknya dipakai untuk lompatan, dan mengapa?
  \item Jelaskan mengapa keadaan permainan disimpan pada kode, bukan pada
        tampilan antarmukanya.
\end{enumerate}

\begin{table}[htbp]
\centering
\caption{Kunci soal konsep Bab XIV beserta indikator yang diukur}
\label{tab:kunci-konsep-ch14}
{\setlength{\tabcolsep}{4pt}\footnotesize%
\begin{tabularx}{\textwidth}{@{}C{0.8cm} C{2.1cm} Y@{}}
\toprule
\headrow \textbf{No} & \textbf{Indikator} & \textbf{Pokok jawaban} \\
\midrule
1 & \ind{a} &
Tingginya $192 \div 100 = 1{,}92$ satuan. Agar tingginya 2 satuan, nilai yang
diperlukan $192 \div 2 = 96$. \\
\addlinespace
2 & \ind{a} &
Karena Unity akan mengolah ulang berkas itu, dan hasilnya untuk seluruh benda
yang memakai berkas tersebut ikut berubah. Ukuran gambar sebaiknya diubah
melalui \istilah{Pixels Per Unit}, agar perubahannya terkendali. \\
\addlinespace
3 & \ind{b} &
Pertama, nama parameter pada kode berbeda dari nama pada jendela
\istilah{Animator}, termasuk perbedaan besar-kecil huruf. Kedua, transisinya
belum digambar, atau syaratnya tidak pernah terpenuhi. \\
\addlinespace
4 & \ind{b} &
Pemicu, sebab lompatan adalah kejadian yang terjadi sekali dan tidak memiliki
lama. Mesin mengembalikan nilai pemicu setelah dipakai, sehingga kode tidak
perlu mengembalikannya sendiri. \\
\addlinespace
5 & \ind{c} &
Agar keadaan permainannya tetap menjadi satu sumber kebenaran, dan tidak
bergantung pada tampilannya. Apabila keadaan disimpan pada tampilan, keadaannya
sulit disimpan, dipulihkan, atau diperiksa terpisah dari antarmukanya. \\
\bottomrule
\end{tabularx}}
\end{table}
